% ID: SEL-D0071-OA-E01
% ESTADO: validada. S071.
Con $0\le y\le5/4$, las cotas de $x$ son $(3-y)/6\le x\le(2-y)/2$. La inferior es positiva y no supera a la superior en ese intervalo. Para $y=0$ salen $(1/2,0)$ y $(1,0)$. Para $y=5/4$ salen $(7/24,5/4)$ y $(3/8,5/4)$.
\begin{center}
\setlength{\unitlength}{0.9mm}
\begin{picture}(97,62)(-8,-8)
\qbezier(0,0)(40,0)(80,0)
\qbezier(0,0)(0,24)(0,48)
\qbezier(35,0)(52.5,0)(70,0)
\qbezier(70,0)(48.125,20)(26.25,40)
\qbezier(26.25,40)(23.333,40)(20.417,40)
\qbezier(20.417,40)(27.7085,20)(35,0)
\put(81,0){$x$}\put(2,48){$y$}
\put(30,-6){\scriptsize $(1/2,0)$}\put(65,-6){\scriptsize $(1,0)$}
\put(4,42){\scriptsize $(7/24,5/4)$}\put(28,40){\scriptsize $(3/8,5/4)$}
\end{picture}
\end{center}
Una función lineal alcanza sus extremos en vértices. En el orden anterior, $F$ vale $1/2$, $1$, $-23/24$ y $-7/8$. El máximo es $\boxed{1}$ sólo en $(1,0)$ y el mínimo es $\boxed{-23/24}$ sólo en $(7/24,5/4)$.
